\documentclass[11pt,a4paper]{article}

\usepackage[utf8]{inputenc}
\usepackage[T1]{fontenc}
\usepackage{geometry}
\usepackage{hyperref}
\usepackage{enumitem}
\usepackage{xcolor}
\usepackage{fancyhdr}
\usepackage{amsmath}
\usepackage{booktabs}
\usepackage{listings}

\geometry{margin=1in}
\hypersetup{
  colorlinks=true,
  linkcolor=blue,
  urlcolor=blue
}

\setlist{itemsep=0.35em, topsep=0.35em}

\pagestyle{fancy}
\fancyhf{}
\lhead{Object Oriented Programming (B.Tech CSE)}
\rhead{Unit 01: Diagnostic Assessment}
\cfoot{\thepage}
\setlength{\headheight}{14pt}

\lstset{
  language=Java,
  basicstyle=\ttfamily\small,
  keywordstyle=\color{blue},
  commentstyle=\color{gray},
  stringstyle=\color{teal},
  showstringspaces=false,
  columns=fullflexible,
  frame=single
}

% Marks per question (total = 100)
\newcommand{\EMarks}{\textbf{[2 marks]} }
\newcommand{\MMarks}{\textbf{[3 marks]} }
\newcommand{\HMarks}{\textbf{[5 marks]} }

\begin{document}

\begin{center}
  {\LARGE \textbf{Diagnostic Assessment -- Unit 01}}\\[0.25em]
  {\large Object Oriented Programming (CSEG1044)}\\[0.25em]
  \normalsize (Java fundamentals + OOP basics; designed to identify learning gaps)
\end{center}

\noindent
\textbf{Instructor:} Tofik Ali \hfill \textbf{Time Limit:} None\\
\textbf{Submission Deadline:} March 31, 2026\\
\textbf{Student Name:} \_\_\_\_\_\_\_\_\_\_\_\_\_\_\_\_\_ \hfill \textbf{Roll No.:} \_\_\_\_\_\_\_\_\\

\vspace{0.6em}
\hrule
\vspace{0.8em}

\textbf{Instructions}
\begin{itemize}[leftmargin=*]
  \item \textbf{This is a diagnostic assessment, not a quiz/test.} The goal is to identify learning gaps so you know what to revise next.
  \item \textbf{Total marks: 100.} Easy: 10 $\times$ 2 = 20, Medium: 10 $\times$ 3 = 30, Hard: 10 $\times$ 5 = 50.
  \item \textbf{No time limit.} Submit on or before \textbf{March 31, 2026}.
  \item \textbf{Important (70\% rule): To have your assignment marks counted, you must score at least 70\% (70/100 or more) in this assessment.}
    If you score below 70\%, you may reattempt and resubmit. Your latest submission will be considered.
  \item \textbf{Marking policy (honest attempt):} Marks are deducted only for (i) not attempting a question and (ii) high similarity with other submissions (copying).
    Correctness is \textbf{not} the main focus; your reasoning and effort matter more.
  \item Answer \textbf{all} questions. For code questions, you may write code snippets (they need not be full programs unless asked).
  \item For MCQs, follow the instruction in the question (some are \textbf{select all that apply}).
\end{itemize}

\vspace{0.6em}

\section*{Easy (10 Questions)}
\begin{enumerate}[label=\textbf{E\arabic*.}, leftmargin=*]
  \item \EMarks \textbf{[MCQ -- Select all that apply]} Which are core OOP principles?
    \begin{enumerate}[label=(\Alph*), leftmargin=*]
      \item Abstraction
      \item Encapsulation
      \item Compilation
      \item Inheritance
      \item Polymorphism
      \item Pagination
    \end{enumerate}

  \item \EMarks \textbf{[MCQ]} If a class is declared \texttt{public class Hello}, what must the file name be?
    \begin{enumerate}[label=(\Alph*), leftmargin=*]
      \item \texttt{Main.java}
      \item \texttt{Hello.java}
      \item Any name is fine
      \item \texttt{hello.Java}
    \end{enumerate}

  \item \EMarks \textbf{[MCQ -- Select all that apply]} Which component contains \texttt{javac}?
    \begin{enumerate}[label=(\Alph*), leftmargin=*]
      \item JDK
      \item JRE
      \item JVM
    \end{enumerate}

  \item \EMarks \textbf{[Subjective]} What is bytecode? Write 2--3 lines.

  \item \EMarks \textbf{[Subjective]} What is the difference between \texttt{==} and \texttt{equals()} for \texttt{String}?
    (3--5 lines)

  \item \EMarks \textbf{[Output]} Predict output:
    \begin{lstlisting}
System.out.println(5/2);
System.out.println(5/2.0);
    \end{lstlisting}

  \item \EMarks \textbf{[Subjective]} Array length is 10. What is the last valid index? Why?

  \item \EMarks \textbf{[Subjective]} What is the type of \texttt{args[i]} inside \texttt{main(String[] args)}?

  \item \EMarks \textbf{[Subjective]} Why can missing \texttt{break} in a \texttt{switch} be dangerous? (2--3 lines)

  \item \EMarks \textbf{[Subjective]} Give one example invariant for a \texttt{BankAccount} or \texttt{Student}.
\end{enumerate}

\section*{Medium (10 Questions)}
\begin{enumerate}[label=\textbf{M\arabic*.}, leftmargin=*]
  \item \MMarks \textbf{[Code]} Write a loop to compute factorial of \texttt{n} (assume \texttt{n >= 0}).

  \item \MMarks \textbf{[Concept + Fix]} You read an \texttt{int} using \texttt{Scanner.nextInt()} and then immediately read a line using
    \texttt{Scanner.nextLine()}, but the line becomes empty. Why? How do you fix it? (4--6 lines)

  \item \MMarks \textbf{[Code]} Write code to read \texttt{n} integers into an array and compute the average.

  \item \MMarks \textbf{[Code]} Write a method \texttt{max(int[] a)} that returns the maximum element.

  \item \MMarks \textbf{[Code]} Write an if--else ladder to convert marks (0--100) into grade:
    \texttt{A (>=80)}, \texttt{B (>=60)}, \texttt{C (>=40)}, else \texttt{F}.

  \item \MMarks \textbf{[Output]} Predict output:
    \begin{lstlisting}
int x = 2 + 3 * 4;
int y = (2 + 3) * 4;
System.out.println(x + " " + y);
    \end{lstlisting}

  \item \MMarks \textbf{[Code]} Write code to reverse a string without using built-in reverse methods.

  \item \MMarks \textbf{[Code]} Given a 2D array \texttt{int[][] m}, write code to print row-wise sums.

  \item \MMarks \textbf{[Design]} Problem: ``Campus resource booking: resources can be booked for a time slot.''\\
    List 4 classes and 1 responsibility per class.

  \item \MMarks \textbf{[Code]} Write a loop that prints numbers 1 to \texttt{n} but skips multiples of 3 (use \texttt{continue}).
\end{enumerate}

\section*{Hard (10 Questions)}
\begin{enumerate}[label=\textbf{H\arabic*.}, leftmargin=*]
  \item \HMarks \textbf{[Program]} Write a program that takes integer command line args and prints:
    count, sum, min, max. If no args are provided, print usage.

  \item \HMarks \textbf{[Program]} Write a palindrome checker that ignores spaces and case. Show a sample run.

  \item \HMarks \textbf{[Program]} Read \texttt{n} marks into an array and print:
    min, max, average, and number of students who scored $\ge$ 40.

  \item \HMarks \textbf{[Program]} Print the first \texttt{n} Fibonacci numbers using a loop.

  \item \HMarks \textbf{[Program]} Implement bubble sort for an integer array and print the sorted result.

  \item \HMarks \textbf{[Explain]} Explain in 8--12 lines: why \texttt{StringBuilder} is preferred over \texttt{String}
    concatenation inside loops.

  \item \HMarks \textbf{[Program]} Count frequency of each character in a string and print only characters that appear (ASCII is fine).

  \item \HMarks \textbf{[Program]} Matrix addition: read two same-size matrices and print their sum.

  \item \HMarks \textbf{[Design + Justify]} For two pairs, choose is-a or has-a and justify (2--3 lines each):
    \begin{enumerate}[label=(\alph*), leftmargin=*]
      \item Car -- Engine
      \item Dog -- Animal
    \end{enumerate}

  \item \HMarks \textbf{[Program]} Write a menu-driven calculator (repeat until exit). Handle divide-by-zero using if--else.
\end{enumerate}

\vfill
\noindent\textit{End of Assessment}

\end{document}

